\documentclass[11pt,a4paper]{article}
\usepackage[T1]{fontenc}\usepackage[utf8]{inputenc}\usepackage{lmodern}
\usepackage[a4paper,margin=2.2cm]{geometry}
\usepackage{amsmath,amssymb}\usepackage{enumitem}\usepackage{booktabs,longtable,array}
\usepackage[hidelinks]{hyperref}
\setlength{\parindent}{0pt}\setlength{\parskip}{0.4em}
\title{\textbf{OP-CONST T1a: the inheritance map $Z_B^{(0)}\!\leftarrow\!
K_\theta^{\rm eff}$ and the fate of $u$-dependence in $q^2/Z_A$}\\\large Recon v1
--- proposal-only; no preprint edits}
\author{Per accepted T1 (rev.~1.1) directive}
\date{12 June 2026}
\begin{document}\maketitle

\section*{0. Scope and firewall}
T1a specifies the inheritance map by which the gauge-kinetic coefficient acquires
its scale from the phase-sector stiffness, and classifies --- mechanism by
mechanism --- whether the physical combination $q^2/Z_A$ retains or cancels
stiffness/background dependence after canonical gauge-field normalisation. This is
a classification-and-specification recon on top of the recorded structure
(\S sec:u1\_emergence Layers 1--3, Appendix A.3, eq:GN\_S0\_common), not a
derivation of $Z_B^{(0)}$; every derivable piece is labelled as such, every
mechanism statement is conditional. Proposal-only; no D1; no preprint edits.

\section{Bookkeeping lemma (derivable now)}
In the $Z$-frame, $\mathcal L=-\tfrac{Z_A}{4}F_{AB}F^{AB}$ with integer charges
$q$ on the holonomy lattice; in the canonical frame, $\hat A=\sqrt{Z_A}\,A$ and
$e=q/\sqrt{Z_A}$. The observable $\alpha_{\rm fs}=e^2/4\pi=q^2/(4\pi Z_A)$ is
invariant under any field rescaling: $u$-dependence cannot be created or removed
by redefinition of $A$. Since the lattice is rigid (T1, V1), the tracking exponent
is fixed exactly:
\[
\zeta_\alpha\;=\;-\,\zeta_{Z_A}.
\]
The inheritance question is therefore well-posed only at the level of the
gauge-invariant pair (charge lattice, $Z_A$), and everything reduces to the
composition of $Z_A$.

\section{Consistency constraint on the inheritance map (conditional lemma)}
The record contains both (i)~the phase-stiffness term
$\tfrac{K_\theta^{\rm eff}}{2}(\partial\theta)^2$ (L$\sim$3750) with its
gauge-covariant completion $\mathcal D_A\theta=\partial_A\theta-eA_A$ used in the
holonomy theorem (eq:gauge\_inv\_phase\_gradient), and (ii)~the massless Maxwell
term as the unique leading kinetic structure (eq:Maxwell\_action).
\textbf{Conditional lemma:} \emph{if} the gauged completion of the recorded
stiffness term is minimal, $\tfrac{K_\theta^{\rm eff}}{2}(\mathcal D_A\theta)^2$,
then it is a gauge-invariant St\"uckelberg mass term: in unitary gauge it yields
$m_A^2=e^2K_\theta^{\rm eff}\sim e^2\phi_0^2$ --- a Planck-scale photon mass,
incompatible with the recorded massless leading-order structure. \textbf{
Consequence:} the inheritance $Z_B^{(0)}\!\leftarrow\!K_\theta^{\rm eff}$
\emph{cannot} be the naive multiplicative transfer through minimal coupling of the
same mode. Either (a)~the combination eaten at the condensate scale is a heavy
gauge direction while the surviving massless field is the orthogonal (unbroken)
combination --- the electroweak-like pattern, consistent with the two-sector
$(Z_B,Z_W,\theta_W)$ structure of Appendix A.3 --- or (b)~the transfer proceeds
through a loop/duality channel rather than tree-level coupling. This single
observation already constrains what ``inherits from'' can mean.

\section{Admissible generation mechanisms and their compositions}
\textbf{M1 (tree-level / projection).} The massless field carries a bare or
projected kinetic coefficient at the condensate scale. Dimensional analysis is
decisive here: $Z$ is dimensionless while $K_\theta^{\rm eff}$ has mass dimension
two, so any tree-level inheritance is \emph{dimensionally forced} to the form
\[
Z\;\sim\;\frac{K_\theta^{\rm eff}}{M^2}
\]
for some condensate mass scale $M$. If $M^2$ is the condensate gap
$m_\varphi^2$, the recorded identity gives
\[
Z\;\sim\;\frac{K_\theta^{\rm eff}}{m_\varphi^2}\;=\;\frac{2S_0^{\rm EFT}}
{c_\perp}\;=\;\frac{1}{2\lambda}\quad\text{(simplest realisation)},
\]
i.e.\ the $S_0$ lever reappears with dimensional necessity: the tree-level
composition is \emph{automatically} an $S_0$-type invariant, parameter-only under
the candidate lemma of T1 \S2(b). Benchmark-level consistency check (no
derivation claimed): an $O(1)\times S_0$-type coefficient is of the same moderate
order as the target $Z_A^{\rm IR}\approx10.9$ (``order ten'', A.1). The dangerous
variant of M1 requires an \emph{external} fixed scale in the denominator
($Z\sim u^2/M_{\rm ext}^2$), which the recorded condensate-only architecture does
not provide.

\textbf{M2 (matter loops).} Integrating out gapped charged excitations generates
$\delta Z_A\sim\sum_f q_f^2\,b_f\ln(\Lambda/m_f)$-type contributions: the
composition is the \emph{charged spectrum} (charges, multiplicities, masses) plus
endpoints. Anchoring class: that of the spectrum and of the running endpoints ---
exactly the C2/C4 territory of T1 and the endpoint-log channel of Step 3
(stress level $\times70$--$470$ if endpoints track). M2 adds no new exposure
beyond what Steps 1--3 already quantify, but it also cannot by itself protect
$\alpha$: it transfers the question to spectrum anchoring ($v_2$-chain,
$\mathcal I_{\rm EW}$).

\textbf{M3 (phase-sector duality / vortex channel).} In the compact-phase sector,
the Maxwell coefficient can arise through the dual description, with composition
set by vortex-core data (core action, core scale) together with
$K_\theta^{\rm eff}$. Core actions are $S_0$-type dimensionless invariants of the
coherent sector; if the recorded structural universality of $S_0$ extends to the
vortex core action, M3 lands in the same protected class as M1. Decisive
computation: the dual map for the compact phase in the ordered branch and the
core-action universality statement.

\section{T1a synthesis}
\begin{enumerate}[leftmargin=1.9em,itemsep=2pt,label=\textbf{S\arabic*.}]
\item In \emph{all three} admissible mechanisms, the natural composition routes
the would-be $u$-dependence either into $S_0$-type invariants (M1 by dimensional
necessity; M3 via core actions) or into the spectrum/endpoint chain already
covered by conditions C2/C4 (M2). An uncancelled power-law tracking
$Z_A\propto u^2$ requires an external fixed denominator scale absent from the
recorded architecture.
\item Condition C1 of T1 is thereby \emph{sharpened and partially closed}: the
tree-level branch of the inheritance is $S_0$-protected by dimensional analysis
(conditional on the T1 candidate lemma); the residual exposure of the
$\alpha$-chain localises in (a)~the endpoint/spectrum logs (quantified at Step 3:
two orders above the anchor if tracking, hence still requiring C4) and (b)~the
OP-grad anchoring of the evaluation scale (C3).
\item The St\"uckelberg consistency lemma (\S2) is a new structural statement in
its own right: it excludes the naive multiplicative inheritance and forces the
electroweak-like projection pattern or a loop/duality channel --- a concrete
constraint on the future OP-SM derivation of $Z_B^{(0)}$.
\item Nothing here is a completed protection theorem. The chain is: T1 candidate
lemma (S$_0$-universality) $+$ T1a dimensional forcing (tree branch) $+$ C3/C4
(OP-grad, endpoints) $\Rightarrow$ Outcome A. The first two links are now
recon-established arguments; the last two remain the open core.
\end{enumerate}

\section{Next targets (proposal-layer, review-gated)}
(T1a$'$)~Select among M1--M3 by deriving which mechanism the recorded emergence
construction actually realises (OP-SM-level work; the St\"uckelberg lemma already
excludes the naive fourth option). (T1b)~OP-grad anchoring of $u_0$ (C3) --- now
the single largest open link of the $\alpha$-chain. (T1c)~Promote the T1
candidate lemma. (T1d)~$\mu$-side selection principle unchanged. No D1/preprint
contact before review of this document.
\end{document}
